% ID: SEL-D0003-OA-E01
% ESTADO: cotejada. Fuente y fragmento cotejados de forma independiente; salida compilada revisada.
% Fuente: Selectividad/sel_2019_matematicas_aplicadas/Suplente Junio/Exámenes Matemáticas Aplicadas CCSS - A y B.pdf
% SHA-256: a59474e6895bd41034da53b7060969cbe8cbf7f8272f39fbf16821c95f06c84b
% Página: 1; opción A; ejercicio 1.
% Encabezado original: EJERCICIO 1. Solo enunciado; no contiene solución.

Consideremos el recinto definido por las siguientes desigualdades:
\[
  7y\leq 15+3x \qquad y\geq x-3 \qquad 3y\geq -x+11
\]

\begin{enumerate}[label=\alph*)]
  \item \textbf{(2 puntos)} Represente gráficamente el recinto anterior y
  calcule sus vértices.
  \item \textbf{(0.5 puntos)} Calcule en qué puntos se alcanzan los valores
  máximo y mínimo de la función $H(x,y)=4x-y-16$ restringida al anterior
  recinto y obtenga dichos valores.
\end{enumerate}
